\section{A multiple-polylogarithm generator for two harmonic tails}
\label{sec:polylog-generator}

The moment reconstruction has a natural analytic generating function.
We use the nested-sum convention
\begin{equation}
 \Li_{r,s}(x,y)=
 \sum_{m>n\ge1}\frac{x^m y^n}{m^r n^s},
 \label{eq:two-variable-mpl}
\end{equation}
initially in its domain of absolute convergence.  This is a depth-two
multiple polylogarithm in two arguments; the notation is chosen to make
the larger and smaller summation indices visible~\cite{Goncharov,BorweinMPL}.

\begin{theorem}[Two-tail polylogarithmic generator]
\label{thm:two-tail-generator}
Let $p,q>1$ be integers and $|z|<1$.  Then
\begin{equation}
 \boxed{\displaystyle
 \calF_{p,q}(z):=\sum_{n=0}^{\infty}A_p(n)A_q(n)z^n
 =\frac{\zeta(p)\zeta(q)
 -\Li_{p,q}(1,z)-\Li_{q,p}(1,z)-\Li_{p+q}(z)}{1-z}.}
 \label{eq:two-tail-generator}
\end{equation}
The same function has the Gamma-normalized double Mellin integral
\begin{equation}
 \calF_{p,q}(z)=
 \frac1{\Gamma(p)\Gamma(q)}
 \int_0^\infty\!\int_0^\infty
 \frac{t^{p-1}u^{q-1}e^{-t-u}}
 {(1-e^{-t})(1-e^{-u})(1-ze^{-t-u})}
 \dd t\dd u.
 \label{eq:two-tail-mellin}
\end{equation}
\end{theorem}

\begin{proof}
Expand the two tails and interchange absolutely convergent sums:
\begin{align*}
 \calF_{p,q}(z)
 &=\sum_{m,k\ge1}\frac1{m^pk^q}
   \sum_{0\le n<\min(m,k)}z^n\\
 &=\frac1{1-z}\sum_{m,k\ge1}
   \frac{1-z^{\min(m,k)}}{m^pk^q}.
\end{align*}
The part without $z^{\min(m,k)}$ is $\zeta(p)\zeta(q)$.  Splitting the
remaining double sum into $m>k$, $k>m$, and $m=k$ gives respectively
$\Li_{p,q}(1,z)$, $\Li_{q,p}(1,z)$, and $\Li_{p+q}(z)$, proving
\eqref{eq:two-tail-generator}.

For \eqref{eq:two-tail-mellin}, use
\[
 A_p(n)=\frac1{\Gamma(p)}\int_0^\infty
 \frac{t^{p-1}e^{-(n+1)t}}{1-e^{-t}}\dd t
\]
and the analogous formula for $A_q(n)$.  Summing the resulting geometric
series in $ze^{-t-u}$ gives the stated integral.  Positivity for real
$z\in[0,1)$ and absolute convergence for $|z|<1$ justify the
interchanges; the complex identity then follows by holomorphy.
\end{proof}

\subsection{Centered weights as one differential operator}

Let
\begin{equation}
 \Eop=z\frac{\dd}{\dd z}+\frac12.
 \label{eq:centered-Euler}
\end{equation}
Because $\Eop z^n=(n+1/2)z^n$, every polynomial weight is encoded by
functional calculus in this single operator:
\begin{equation}
 P(\Eop)\calF_{p,q}(z)
 =\sum_{n=0}^{\infty}P(n+\tfrac12)A_p(n)A_q(n)z^n.
 \label{eq:operator-tail}
\end{equation}
Likewise,
\begin{equation}
 Q(\Eop)\frac1{1-z}
 =\sum_{n=0}^{\infty}Q(n+\tfrac12)z^n.
 \label{eq:operator-polynomial}
\end{equation}
This gives a compact polylogarithmic normal form for every identity in
\cref{thm:ordinary-reconstruction}.

\begin{corollary}[Polylogarithmic boundary reconstruction]
\label{cor:polylog-boundary}
For $b\ge6$ and $0\le k<b$,
\begin{equation}
 \boxed{\displaystyle
 \lim_{z\uparrow1}
 \left\{
 P_{b,k}(\Eop)\calF_{10,2b}(z)
 -Q_{b,k}(\Eop)\frac1{1-z}
 \right\}=\zeta(2k),}
 \label{eq:polylog-boundary}
\end{equation}
where $\calF_{10,2b}$ is given explicitly by
\eqref{eq:two-tail-generator}.  The limit may be taken radially in the
unit disk.
\end{corollary}

\begin{proof}
Equations \eqref{eq:operator-tail} and
\eqref{eq:operator-polynomial} show that the expression in braces has
Taylor coefficients equal to the summand of
\eqref{eq:ordinary-reconstruction}.  That coefficient series is
absolutely summable, so Abel's theorem identifies the radial limit with
its sum.
\end{proof}

Repeated argument differentiation introduces no new function class:
for every integer $j\ge0$,
\begin{equation}
 \left(z\frac{\dd}{\dd z}\right)^j\Li_{r,s}(1,z)
   =\Li_{r,s-j}(1,z),
 \qquad
 \left(z\frac{\dd}{\dd z}\right)^j\Li_r(z)=\Li_{r-j}(z).
 \label{eq:order-lowering}
\end{equation}
Thus \eqref{eq:polylog-boundary} is a finite identity among ordinary
polylogarithms, depth-two multiple polylogarithms with shifted orders,
and rational functions.  Formula \eqref{eq:two-tail-mellin} supplies a
parallel Gamma-integral representation of the same boundary values.

\begin{remark}[Why the numerator is a divided tail]
The factor $(1-z)^{-1}$ in \eqref{eq:two-tail-generator} is not an
artifact of summation.  It records the finite geometric prefix
$1+z+\cdots+z^{\min(m,k)-1}$.  The numerator therefore measures exactly
the difference between the unrestricted product $\zeta(p)\zeta(q)$ and
the three regions in which the smaller index has already crossed the
prefix boundary.
\end{remark}
